\section{Pointwise deconcentration at endpoint decision boundaries}
\label{sec:v47-endpoint-decisions}
A section-decision boundary is not a singularity of a fixed physical
collision word. A roof flow may cross such a boundary while retaining
the same displacement and collision count. The return index can then
change, but by no more than the number of decisions allowed to change.
We use this finite enlargement only in a positive upper comparison.
The density being estimated is always that of the original exact
$(n,k,m)$ event. This distinction permits a pointwise estimate for part
of the collapsing-margin source, without a positive lower bound for
its decision distances.

Write $c=\nu(Y_R^*)$, $q=47/53$, and use $e_j,e'_j$ and $\Theta$ from
Sections~\ref{sec:v45-critical-collars} and
\ref{sec:v46-window-transport}. Fix an integer $J\ge0$ and first take
$m>2J+2$. Put
\begin{align}
 \mathcal E_{m,J}&=\{0,\ldots,J\}\cup\{m-J,\ldots,m\},\notag\\
 r_J&=2J+2,\qquad
 \mathcal I_{n,J}=\{\ell\in\Z:|\ell-n|\le r_J\}.
                                      \label{eq:v47-free-decisions}
\end{align}
There are at most $r_J$ occupation summands in this set of free
decisions; the terminal decision at $m$ is not an occupation summand.
There is no restriction on the initial or terminal section membership
in an auxiliary collision comparison below. Its measure is $\nu/c$,
which agrees with $\nu_R^*$ on the original initial section, but is
not a probability on the whole cylinder.

For $0<\eta<\varepsilon_0$ define
\begin{equation}\label{eq:v47-partial-guard}
 G_{m,R}^{\eta,J}=
 \prod_{j=0}^{m}\Theta\!\left(\frac{c_j}{\eta e_j}\right)
 \prod_{j=0}^{m-1}\Theta\!\left(\frac{\Delta_j}{\eta e'_j}\right)
 \prod_{j\notin\mathcal E_{m,J}}
                    \Theta\!\left(\frac{s_j}{\eta e_j}\right).
\end{equation}
Only the section factors in $\mathcal E_{m,J}$ have been omitted.
All incidence and clearance factors remain. Let $\mathsf P^{\eta,J}$
denote the corresponding set of once-thresholded regular histories.
For a second parameter $0<\varepsilon<\varepsilon_0$ put
\begin{align}
 E_{m,R}^{\varepsilon,J}
  &=\prod_{j\in\mathcal E_{m,J}}
                     \Theta\!\left(\frac{s_j}{\varepsilon e_j}\right),
                                                         \notag\\
 d^{\eta,\varepsilon,J,a}_{n,k,m,R}(t)\,dt
  &=(T_{n,R})_*\bigl(\one_{\Pi_{n,k,m,R}}a
             G_{m,R}^{\eta,J}(1-E_{m,R}^{\varepsilon,J})\nu_R^*\bigr).
                                      \label{eq:v47-end-source}
\end{align}
The insertion may be any measurable function with $|a|\le M$.
Absolute continuity follows by domination by the complete physical
density. The two margin parameters will be set equal only after the
pointwise estimate has been proved.

\subsection{Small endpoint sets in the arithmetic local upper bound}
We require a version of the endpoint estimate without section factors.
Use the uncentered collision occupation $A_{m,R}$, displacement
$K^{\rm c}_{m,R}$, and roof $S_{m,R}$, so that on the original event
\[
 (K^{\rm c}_{m,R},A_{m,R},S_{m,R})=(k,n,T_{n,R}).
\]
In particular $A_{m,R}$ sums the visits at times $0,\ldots,m-1$.

\begin{lemma}[Small endpoint mass and an exact collision coefficient]
\label{lem:v47-small-endpoint-local}
Let $0\le a_R,d_R\le1$ be smooth endpoint functions with common
collision-norm bounds and parameter continuity, for the fixed choices
under consideration. Suppose $\nu(a_R)\le\alpha$ and
$\nu(d_R)\le\beta$. For every fixed bounded interval $I$,
\begin{align}
 &\limsup_{m\to\infty}\sup_{R,k,\ell,t}m^2
 \int a_R(x)d_R(T_R^m x)\one_{\{K^{\rm c}_{m,R}=k,\ A_{m,R}=\ell,\
                                      S_{m,R}-t\in I\}}\,d\nu(x)
                 \le C\alpha\beta |I|.
                                      \label{eq:v47-small-endpoint-local}
\end{align}
The constant does not depend on their fixed smooth norms. The
threshold in the collision limit may depend on those norms and on
$I$. The same conclusion follows by positive smooth upper
approximation of endpoint sets.
\end{lemma}
\begin{proof}
Apply the exact full-occupation pairing and the finite-branch argument
of Theorem~\ref{thm:v41-uniform-transition}, with endpoints $a_R,d_R$
in place of $\eta_R,\eta_R$. The eigenvalues, moving maxima and
quadratic matrices are unchanged. On every branch surviving along a
convergent radius sequence, the limiting physical projection is an
actual resonance. Its scalar amplitude has modulus
\[
 |\nu(d_R q_\gamma)\nu(a_R\overline q_\gamma)|
                                     \le\alpha\beta,
\]
since $|q_\gamma|=1$. The complex Gaussian integrals have a common
upper bound for every real target, and the finite number of surviving
peak charts is independent of the endpoint functions and the outer
roof band. Every limiting roof resonance has roof coordinate zero.
This is the same subsequence argument used in
Lemma~\ref{lem:v45-two-strip-local}; it does not assume continuity
of a generator of the arithmetic group.

For a positive Fourier majorant of $\one_I$, the limiting roof factor
is its mass $|I|+2\pi/B$. Positivity of the unchanged collision
coefficient gives the upper inequality before taking limits. First
send $m$ to infinity with $B$ and the endpoints fixed, and then send
$B$ to infinity. This proves \eqref{eq:v47-small-endpoint-local}.
An upper approximation is taken with fixed smooth norms before either
limit, and its excess mass is removed afterwards. No rate uniform in
a shrinking endpoint set or roof interval is asserted.
\end{proof}

For a fixed constant $a_*>0$ define the forward and backward endpoint
sets
\begin{align}
 U^+_{R,J}(s)&=\{x:\exists\,0\le j\le J,
       \operatorname{dist}(T_R^jx,\partial Y_R^*)<a_*s q^{j/4}\},\notag\\
 U^-_{R,J}(s)&=\{y:\exists\,0\le j\le J,
       \operatorname{dist}(T_R^{-j}y,\partial Y_R^*)<a_*s q^{j/4}\}.
                                      \label{eq:v47-endpoint-bad-sets}
\end{align}
Singular finite histories may be included in these sets; they are
null and are included in their smooth upper approximations.

\begin{lemma}[Finite-block endpoint envelopes]
\label{lem:v47-endpoint-envelopes}
For fixed $J$ and $s>0$, both sets in
\eqref{eq:v47-endpoint-bad-sets} admit positive smooth upper
approximations, locally in the parameter and uniformly on a finite
parameter cover, with masses at most $C s+\xi$ for arbitrary
$\xi>0$. The constant $C$ is independent of $J$. Consequently
\begin{align}
 &\limsup_{m\to\infty}\sup_{R,k,\ell,t}m^2
 \nu\{K^{\rm c}_{m,R}=k,A_{m,R}=\ell,S_{m,R}-t\in I,\notag\\
 &\hspace{40mm}x\in U^+_{R,J}(s)
                       \text{ or }T_R^mx\in U^-_{R,J}(s)\}
                                      \le C s|I|.
                                      \label{eq:v47-bad-end-local}
\end{align}
Also the same collision event restricted by
$|p_0|,|p_m|\le d$ has normalized limsup at most $Cd^2|I|$.
\end{lemma}
\begin{proof}
Invariance and the section-neighborhood estimate give
\[
 \nu(U^\pm_{R,J}(s))
       \le C s\sum_{j=0}^Jq^{j/4}\le C s.
\]
We explain why the endpoint regularity needed here follows from
finite-time geometry rather than from a presumed long-pullback norm.
For fixed $J$ the finite union of forward or backward collision
singularities and their preimages is null. At a fixed limiting radius,
remove a closed neighborhood of this set and of grazing with mass
less than $\xi$. On its compact complement every finite itinerary
and every distance function in \eqref{eq:v47-endpoint-bad-sets}
persists uniformly under a small parameter change. A closed bad set
with threshold $a_*s$ is then contained in the corresponding open
set with threshold $2a_*s$, together with that neighborhood.
The latter union has mass at most $Cs+C\xi$, again by invariance.
A smooth function equal to one on the compact bad set and on a
slightly smaller singular neighborhood can be supported in a slightly
larger such union. Its mass is at most $Cs+C\xi$. This also covers
states close to the boundary of the compact complement. Finite
parameter compactness and a subordinate continuous partition in $R$
give common smooth norms and the required scalar continuity; the
smooth norms may depend on $J,s,\xi$. Rescale $\xi$ to obtain the
stated excess. This argument uses a doubled distance threshold and
therefore does not require a regular value for that distance.

Apply Lemma~\ref{lem:v47-small-endpoint-local} separately to the
initial and terminal upper envelopes, with the other endpoint one.
The union is bounded by their sum. Remove $\xi$ after the collision
and band limits to prove \eqref{eq:v47-bad-end-local}. For the last
assertion use smooth strip functions of mass at most $2d$ at both
ends, exactly as in Lemma~\ref{lem:v45-two-strip-local}, but without
multiplying by the section indicators.
\end{proof}

\subsection{Roof transport through finitely many section cuts}
\begin{lemma}[Continuation with free endpoint decisions]
\label{lem:v47-free-decision-continuation}
Every state in $\mathsf P^{\eta,J}$ has the endpoint continuation
of Lemma~\ref{lem:v46-regular-endpoints}, of radius $r\eta^3$.
It preserves the center word, $K^{\rm c}_{m,R}$, $m$, and all section
decisions outside $\mathcal E_{m,J}$. On a continuation starting with
$A_{m,R}=n$ the occupation always belongs to $\mathcal I_{n,J}$.
The Hessian and logarithmic-density bounds hold with $\eta$ in place
of $\varepsilon$. Moreover
\begin{equation}\label{eq:v47-scaled-decision-speed}
 \max_{0\le j\le m}e_j^{-1}|\nabla s_j|\le C\eta^{-1}
                         \quad\text{almost everywhere}.
\end{equation}
The physical density is $w\,du\,dv=|\partial_u\partial_vF|
(4\pi Rc)^{-1}du\,dv$, including when an endpoint crosses the section.
\end{lemma}
\begin{proof}
The endpoint response and density distortion use incidence and
clearance, not section membership. The continuation proof checks a
section distance only to preserve that particular decision. Omit those
checks at the free indices and retain every other half-margin check.
This proves the same endpoint square and all physical estimates.
The collision-state response bound is
$C\eta^{-1}q^{3\min(j,m-j)/4}$, up to a fixed adjacent-index factor.
Distances to the finite rectangle boundary are Lipschitz, including
at its corners. Division by $e_j$ proves
\eqref{eq:v47-scaled-decision-speed} without a factor depending on
$J$ or $m$. A count of the free summands of $A_m$ gives
$|A_m-n|\le r_J$. The terminal visit is treated as a membership
condition only and is dropped in the upper comparison.

A section cut does not change the reduced physical action or its
cross Hessian. Thus these functions glue across a free decision cut.
The displayed density is precisely $\nu/c$ on the continued word;
no extension of $\nu_R^*$ as a probability or new section
normalization is used.
\end{proof}

\begin{proposition}[A pointwise bound off the normal endpoints]
\label{prop:v47-high-gradient-boundary}
Let $A$ be any measurable subset of
$\Pi_{n,k,m,R}\cap\mathsf P^{\eta,J}$ with
$|\nabla F|\ge\delta$ and
$s_j<2\varepsilon e_j$ at some $j\in\mathcal E_{m,J}$.
Its roof density, for the original source $\nu_R^*$, satisfies
\begin{equation}\label{eq:v47-high-gradient-boundary}
 \limsup_{m\to\infty}\sup_{R,n,k,A}m^2\|\rho_A\|_\infty
                                      \le C(J+1)\varepsilon.
\end{equation}
Here $J\ge0$, $\eta,\varepsilon>0$ and $0<\delta<1/10$ are fixed
before the limit, with $J$ an integer. The constant $C$ is independent of these choices.
\end{proposition}
\begin{proof}
Use $V=\nabla F/|\nabla F|^2$. With a fixed sufficiently small $c_2$,
choose the fixed roof width
\begin{equation}\label{eq:v47-free-flow-width}
 h=c_2\min\{\eta^3\delta^2,\,\varepsilon\eta\delta\}.
\end{equation}
The first term ensures the continuation, the lower gradient
$\delta/2$ and the physical density Jacobian between $e^{-1}$ and $e$,
as in Lemma~\ref{lem:v46-positive-flow-box}. The second term and
\eqref{eq:v47-scaled-decision-speed} ensure that at least the initially
bad free decision still has $s_j<3\varepsilon e_j$ throughout the flow.
The flow is allowed to cross the free section cuts; it crosses no
physical-word or protected middle-decision boundary.

For almost every level $F=t$, the map from that level times $(-h,h)$
is injective. The roof identifies the flow time and uniqueness
identifies the starting point. Across free decision cuts uniqueness
is applied to the single physical-word vector field, not to separate
cut domains. Different center words remain disjoint. Its image lies
in the event
\[
 K^{\rm c}_{m,R}=k,\quad A_{m,R}\in\mathcal I_{n,J},\quad
 |S_{m,R}-t|<h,
\]
with $x\in U^+_{R,J}(\varepsilon)$ or
$T_R^mx\in U^-_{R,J}(\varepsilon)$, taking $a_*=3$.
Weighted coarea consequently gives
\begin{equation}\label{eq:v47-flow-to-relaxed-window}
 \rho_A(t)\le\frac{e}{2hc}
   \sum_{\ell\in\mathcal I_{n,J}}
   \nu\{K^{\rm c}_{m,R}=k,A_{m,R}=\ell,|S_{m,R}-t|<h,
                                     \text{bad endpoint set}\}.
\end{equation}
There are at most $4J+5$ terms. Apply
\eqref{eq:v47-bad-end-local} to this fixed $h$ and cancel the factor
$2h$. This proves \eqref{eq:v47-high-gradient-boundary}.
The cancellation occurs after a fixed-window local upper bound;
it is not a substitution of a shrinking window into that theorem.
The left side retains the prescribed original return index.
\end{proof}

\subsection{Critical completion with the same finite decision allowance}
\begin{lemma}[Low-gradient completion and cluster bound]
\label{lem:v47-free-critical-cluster}
For fixed $\eta,J$ and $0<\delta<d_*\eta^3$, the density of any
positive restriction of
$\Pi_{n,k,m,R}\cap\mathsf P^{\eta,J}\cap\{|\nabla F|\le2\delta\}$
satisfies
\begin{equation}\label{eq:v47-free-low-gradient}
 \limsup_{m\to\infty}\sup_{R,n,k}m^2\|\rho\|_\infty
                                      \le C(J+1)\delta^2.
\end{equation}
The source is still unnormalized. No endpoint section membership is
required at the completed critical center.
\end{lemma}
\begin{proof}
Complete each source point in its free-decision endpoint square to
the unique nearby normal-to-normal critical center. The proof of
Lemma~\ref{lem:v46-near-critical-completion} uses strong convexity
and the preserved physical and middle margins, so applies unchanged
after omitting the free section checks. The center is within
$C\delta$ and its roof within $C\delta^2$. It is
$\eta/2$-protected at every retained margin. Its middle occupation
is the original one; its total occupation can differ from $n$ by
at most $r_J$.

Use the actual mass collars of these centers in the physical word,
with source $\nu/c$. Their intrinsic coefficients are
$J_z=2\pi w_z(z)/\sqrt{\det\nabla^2F_z(z)}$ in this same measure.
For $0<r<h_{\eta/2}$, a collar of roof width $r$ has mass at least
$a_*rJ_z$ and both endpoint tangential velocities at most $A_*\sqrt r$.
Take each physical collar only once even if its free decisions split
it into several return-index pieces. Physical endpoint injectivity
and the inherited collar construction give disjoint original-source
collars. The continuation keeps the middle occupation, so every point
of their union has $A_m\in\mathcal I_{n,J}$. Here the middle count of
a center is the same as that of its starting source: both occupations
lie between that count and that count plus $r_J$. Thus their mutual
difference is at most $r_J$, not an iterated enlargement.

For centers with critical roof in $[t-r,t+r]$, the positive collar
comparison now uses the interval $[t-r,t+2r]$, the two normal endpoint
strips, and at most $4J+5$ occupation coefficients. The last assertion
of Lemma~\ref{lem:v47-endpoint-envelopes} bounds its normalized
limsup by $C(J+1)r^2$. Divide by the collar mass factor $a_*r$ to
obtain the cluster bound $C(J+1)r$ for the sum of intrinsic coefficients.
The collar-density upper estimate then bounds the completed source
by this cluster on a roof interval of width $C\delta^2$.
This proves \eqref{eq:v47-free-low-gradient}. All collar widths and
strip widths are fixed in the collision limit.
\end{proof}

\begin{theorem}[Endpoint decision boundaries are pointwise small]
\label{thm:v47-endpoint-boundary-smallness}
For fixed $J\ge0$, $\eta,\varepsilon\in(0,\varepsilon_0)$,
\begin{equation}\label{eq:v47-endpoint-density-smallness}
 \limsup_{m\to\infty}\sup_{R,n,k,\,|a|\le M}
 m^2\|d^{\eta,\varepsilon,J,a}_{n,k,m,R}\|_\infty
                                      \le CM(J+1)\varepsilon.
\end{equation}
Consequently the unchanged kernel $K_B$ satisfies
\begin{equation}\label{eq:v47-endpoint-all-band}
 \limsup_{m\to\infty}\sup_{B>0}\sup_{R,n,k,\,|a|\le M}
 m^2\|d^{\eta,\varepsilon,J,a}-K_B*d^{\eta,\varepsilon,J,a}\|_\infty
                                      \le CM(J+1)\varepsilon.
\end{equation}
The supremum over insertions includes arbitrary bounded measurable
path functions, by domination. No derivatives of the insertion are
needed. The constants are independent of the fixed protection $\eta$;
the asymptotic threshold is not asserted independent of it.
\end{theorem}
\begin{proof}
It suffices to estimate the positive source with $a=1$. On its support
some free decision has $s_j<2\varepsilon e_j$. Split at
$|\nabla F|=\delta$, applying Proposition~\ref{prop:v47-high-gradient-boundary}
to the high-gradient part and Lemma~\ref{lem:v47-free-critical-cluster}
to the rest. The normalized limsup is at most
$C(J+1)(\varepsilon+\delta^2)$. Send $\delta$ to zero only after that
limsup, with $\eta,\varepsilon,J$ fixed, to prove
\eqref{eq:v47-endpoint-density-smallness}. Radon--Nikodym domination
gives the uniform bound for $|a|\le M$.
Finally $\|K_B\|_1=\|K_1\|_1$ and
$\|f-K_B*f\|_\infty\le(1+\|K_1\|_1)\|f\|_\infty$ prove
\eqref{eq:v47-endpoint-all-band}. The all-band statement is a
convolution consequence of pointwise density control, not a spectral
bound with growing frequency.
\end{proof}

\begin{corollary}[Finite jumps of this boundary component]
\label{cor:v47-endpoint-jumps}
Where finite one-sided traces of the unweighted boundary density
exist, the height of each of its jumps satisfies the bound
\eqref{eq:v47-endpoint-density-smallness}, with twice its constant.
In particular these jump heights are $o(m^{-2})$ in the ordered limits
$m\to\infty$, then $\varepsilon\downarrow0$, for fixed $J,\eta$.
This statement does not bound the sum of the absolute heights of all
jumps and does not apply to the remaining grazing or clearance source.
\end{corollary}
\begin{proof}
Both finite traces of an essentially bounded density have modulus at
most its essential bound whenever they exist. Their difference is at
most twice that bound. Apply Theorem~\ref{thm:v47-endpoint-boundary-smallness}.
The intrinsic right-trace convention is unchanged.
\end{proof}
